% Reviewed translation: PDF pages 23--31 / printed pages 9--17.
% The original scan is authoritative; environment headings remain English.
\MathBookSourcePage{23}

除了边界值算子 $\gamma_0$ 以外, 还需要法向导数的迹 $\gamma_1u$. 对 $\mathcal D(\bar\Omega)$ 中的 $u$, 它定义为

\begin{equation}\label{eq:normal-derivative-trace}
\gamma_1u=\partial u/\partial n=\sum_{i=1}^N\gamma_0(\partial u/\partial x_i)n_i.
\end{equation}

其中 $\mathbf n=(n_1,\ldots,n_N)$ 表示 $\Gamma$ 的单位外法向量. 于是可以如下补充
Theorem~\ref{thm:trace-operator} 的陈述.

\setcounter{statement}{5}
\renewcommand{\theHstatement}{theorem.\thesection.\arabic{statement}}
\begin{Theorem}\label{thm:normal-derivative-trace}
保持 Theorem~\ref{thm:trace-operator} 的假设, 并令 $l\geq1$. 定义在 $\mathcal D(\bar\Omega)$ 上的映射

\[
u\mapsto\{\gamma_0u,\gamma_1u\}
\]

存在唯一的线性连续延拓, 该延拓是从下列空间到下列空间的满射算子:

\[
W^{s,p}(\Omega)\quad\text{onto }W^{s-1/p,p}(\Gamma)\times W^{s-1-1/p,p}(\Gamma).
\]

此外, 在 $W^{2,p}(\Omega)$ 中有如下刻画:

\[
\operatorname{Ker}\gamma_0\cap\operatorname{Ker}\gamma_1=W_0^{2,p}(\Omega).
\]
\end{Theorem}

\setcounter{statement}{0}
\renewcommand{\theHstatement}{remark.\thesection.\arabic{statement}}
\begin{Remark}\label{rem:polygon-boundary-traces}
如果 $\Omega$ 的边界有角点, 它的法向量会发生跳跃, 因而无论 $u$ 多么光滑, $\partial u/\partial n$ 显然都会比较粗糙. 尽管如此, 仍可将 Theorem~\ref{thm:normal-derivative-trace} 的陈述略作推广. 假设 $\Omega$ 是有界二维多边形; 以 $\Gamma_j$ 表示 $\Gamma$ 的各条边, 以 $\mathbf n_j$ 表示相应的单位外法向量, $1\leq j\leq J$. 则映射

\[
u\mapsto(\partial u/\partial n_j;\ 1\leq j\leq J)
\]

是从 $W^{k+2,p}(\Omega)$ 到

\[
\prod_{j=1}^J W^{k+1-1/p,p}(\Gamma_j)
\]

的线性连续满射, 其中 $k\geq0$ 是任意整数, $p\in(1,\infty)$ 是任意实数. 注意, $\partial u/\partial n$ 在 $\Gamma$ 的顶点处没有匹配条件.

$u$ 的边界值的情形稍微复杂一些, 因为在 $\Gamma$ 的顶点处通常存在匹配条件. 为简单起见, 只取 $u\in W^{2,p}(\Omega)$. 以 $S_j$ 表示 $\Gamma$ 的顶点, 并约定 $S_{J+1}=S_1$. 映射

\[
u\mapsto(h_j=u|_{\Gamma_j};\ 1\leq j\leq J)
\]

是从 $W^{2,p}(\Omega)$ 到

\[
\prod_{j=1}^J W^{2-1/p,p}(\Gamma_j)
\]

中由下列相容条件确定的子空间的线性连续满射:

\[
h_j(S_{j+1})=h_{j+1}(S_{j+1}),\qquad1\leq j\leq J.
\]
\end{Remark}

\MathBookSourcePage{24}

最后给出 Green 公式的两个有用应用, 以此结束本节.

\setcounter{statement}{3}
\renewcommand{\theHstatement}{lemma.\thesection.\arabic{statement}}
\begin{Lemma}\label{lem:green-formula-applications}
设 $\Omega$ 是 $\mathbb R^N$ 的有界开子集, 其边界 $\Gamma$ 是 Lipschitz 连续的.

$1^\circ)$ 对 $H^1(\Omega)$ 中的 $u$ 和 $v$, 以及 $1\leq i\leq N$, 有:

\begin{equation}\label{eq:green-first-identity-component}
\int_\Omega u(\partial v/\partial x_i)\,dx=-\int_\Omega(\partial u/\partial x_i)v\,dx+\int_\Gamma\gamma_0(uv)n_i\,ds.
\end{equation}

$2^\circ)$ 如果还有 $u\in H^2(\Omega)$, 则有:

\begin{equation}\label{eq:green-gradient-component-sum}
\sum_{i=1}^N\int_\Omega\frac{\partial u}{\partial x_i}\frac{\partial v}{\partial x_i}\,dx=-\sum_{i=1}^N\int_\Omega\frac{\partial^2u}{\partial x_i^2}v\,dx+\sum_{i=1}^N\int_\Gamma\gamma_0\left(\frac{\partial u}{\partial x_i}v\right)n_i\,ds.
\end{equation}

采用通常的记号

\[
\Delta u=\sum_{i=1}^N\frac{\partial^2u}{\partial x_i^2},\qquad \mathbf{grad}\,u=(\partial u/\partial x_1,\ldots,\partial u/\partial x_N),
\]

式 \eqref{eq:green-gradient-component-sum} 变为

\begin{equation}\label{eq:green-first-identity}
(\mathbf{grad}\,u,\mathbf{grad}\,v)=-(\Delta u,v)+\int_\Gamma(\partial u/\partial n)\gamma_0v\,ds.
\end{equation}
\end{Lemma}

\subsection{抽象椭圆理论}\label{subsec:abstract-elliptic-theory}

本节简要介绍研究椭圆型线性偏微分方程时使用的一种基本工具.

设 $V$ 是实 Hilbert 空间, 其范数记作 $\|\cdot\|_V$; 设 $V'$ 是它的对偶空间, 并以 $\langle\cdot,\cdot\rangle$ 表示 $V'$ 与 $V$ 之间的对偶配对. 设 $(u,v)\mapsto a(u,v)$ 是 $V\times V$ 上的实双线性形式, $l$ 是 $V'$ 中的元素, 考虑如下问题:

\begin{equation}\tag{P}\label{eq:abstract-variational}
\text{Find }u\in V\text{ such that}\qquad a(u,v)=\langle l,v\rangle\quad\forall v\in V.
\end{equation}

下面这个 Theorem 由 Lax \& Milgram [49] 给出.

\setcounter{statement}{6}
\renewcommand{\theHstatement}{theorem.\thesection.\arabic{statement}}
\begin{Theorem}[Lax--Milgram]\label{thm:lax-milgram}
假设 $a$ 在 $V$ 上连续且椭圆, 即存在两个常数 $M$ 和 $\alpha>0$, 使得

\begin{equation}\label{eq:lax-milgram-continuity}
|a(u,v)|\leq M\|u\|_V\|v\|_V\qquad\forall u,v\in V
\end{equation}

以及

\begin{equation}\label{eq:lax-milgram-ellipticity}
a(v,v)\geq\alpha\|v\|_V^2\qquad\forall v\in V.
\end{equation}

则 Problem~\eqref{eq:abstract-variational} 在 $V$ 中存在唯一解 $u$. 此外, 映射 $l\mapsto u$ 是从 $V'$ 到 $V$ 的同构.
\end{Theorem}

\MathBookSourcePage{25}

\setcounter{statement}{1}
\renewcommand{\theHstatement}{corollary.\thesection.\arabic{statement}}
\begin{Corollary}\label{cor:lax-milgram-energy-minimizer}
当 $a$ 对称时, 即 $a(u,v)=a(v,u)$ 对所有 $u,v\in V$ 成立,
Problem~\eqref{eq:abstract-variational} 的解 $u$ 也是 $V$ 中使下列二次泛函(也称为能量泛函)达到最小值的唯一元素:

\begin{equation}\label{eq:energy-functional}
J(v)=(1/2)a(v,v)-\langle l,v\rangle.
\end{equation}
\end{Corollary}

\subsection{Example 1: Laplace 算子的 Dirichlet 问题}\label{subsec:laplace-dirichlet-example}

在所有例子中, 都假设 $\Omega$ 有界且 $\Gamma$ 是 Lipschitz 连续的. 考虑如下非齐次 Dirichlet 问题:

给定 $H^{-1}(\Omega)$ 中的 $f$ 和 $H^{1/2}(\Gamma)$ 中的 $g$, 求函数 $u$, 使得

\noindent Problem (D):
\begin{align}
-\Delta u&=f&&\text{in }\Omega,\label{eq:laplace-dirichlet-pde}\\
u&=g&&\text{on }\Gamma.\label{eq:laplace-dirichlet-boundary}
\end{align}

把这个问题表述为 Problem~\eqref{eq:abstract-variational} 的形式. 令 $V=H_0^1(\Omega)$, 并置

\[
a(u,v)=(\mathbf{grad}\,u,\mathbf{grad}\,v).
\]

显然, $a$ 在 $H_0^1(\Omega)^2$ 上连续, 并且由
Theorem~\ref{thm:poincare-friedrichs},

\[
a(v,v)=\|\mathbf{grad}\,v\|_{0,\Omega}^2=|v|_{1,\Omega}^2\geq C\|v\|_{1,\Omega}^2.
\]

此外, 由于 $H^{1/2}(\Gamma)$ 是 $H^1(\Omega)$ 中 $\gamma_0$ 的值域空间, 取满足 $u_0=g$ on $\Gamma$ 的 $u_0\in H^1(\Omega)$, 并考察如下问题:

\noindent Problem $(D')$: 求 $u\in H^1(\Omega)$, 使得
\begin{align}
u-u_0&\in H_0^1(\Omega),\label{eq:laplace-dirichlet-lifting}\\
a(u-u_0,v)&=\langle f,v\rangle-a(u_0,v)
&&\forall v\in H_0^1(\Omega).\label{eq:laplace-dirichlet-variational}
\end{align}

由于 $a$ 连续, 映射 $v\mapsto\langle f,v\rangle-a(u_0,v)$ 属于 $H^{-1}(\Omega)$.
因此, 由 Lax \& Milgram Theorem, Problem $(D')$ 在 $H^1(\Omega)$ 中存在唯一解 $u$.

只需再证明可以把 $u$ 刻画为 Problem $(D)$ 的唯一解. 在式
\eqref{eq:laplace-dirichlet-variational} 中取 $v\in\mathcal D(\Omega)$, 得到

\[
a(u,v)=-\langle\Delta u,v\rangle=\langle f,v\rangle\qquad\forall v\in\mathcal D(\Omega).
\]

因此 $u$ 满足 \eqref{eq:laplace-dirichlet-lifting} 和 \eqref{eq:laplace-dirichlet-pde}.
反过来, 由 $\mathcal D(\Omega)$ 在 $H_0^1(\Omega)$ 中的稠密性, $(D_1)$ 的每个解都是 $(D')$ 的解. 但

\[
u-u_0\in H_0^1(\Omega)\quad\text{iff }u=g\quad\text{on }\Gamma,
\]

因此 Problems $(D_1)$ 和 $(D)$ 相同.

\MathBookSourcePage{26}

关于 $u$ 的正则性, 由 Lax \& Milgram Theorem 可知, 映射 $l\mapsto u-u_0$ 是从
$H^{-1}(\Omega)$ 到 $H_0^1(\Omega)$ 的同构. 因此,

\[
\|u-u_0\|_{1,\Omega}\leq C_2\|l\|_{-1,\Omega}.
\]

显然,

\[
\|l\|_{-1,\Omega}\leq\|f\|_{-1,\Omega}+\|u_0\|_{1,\Omega}.
\]

从而

\[
\|u\|_{1,\Omega}\leq C_3\{\|f\|_{-1,\Omega}+\|u_0\|_{1,\Omega}\}\quad\forall u_0\in H^1(\Omega)\text{ such that }u_0=g\text{ on }\Gamma.
\]

根据 Definition~\eqref{eq:trace-space-norm}, 这意味着

\[
\|u\|_{1,\Omega}\leq C_3\{\|f\|_{-1,\Omega}+\|g\|_{1/2,\Gamma}\}.
\]

于是证明了如下 Proposition.

\setcounter{statement}{0}
\renewcommand{\theHstatement}{proposition.\thesection.\arabic{statement}}
\begin{Proposition}\label{prop:laplace-dirichlet-wellposedness}
Problem (D) 在 $H^1(\Omega)$ 中存在唯一解 $u$, 并且存在常数 $C=C(\Omega)$, 使得

\begin{equation}\label{eq:dirichlet-stability-h1}
\|u\|_{1,\Omega}\leq C\{\|f\|_{-1,\Omega}+\|g\|_{1/2,\Gamma}\},
\end{equation}

即 $u$ 连续依赖于 Problem (D) 的数据.
\end{Proposition}

当 $f$ 和 $g$ 具有更高正则性时, 自然可以预期 Problem (D) 的解 $u$ 也更光滑. 下一个 Theorem 给出 $u$ 的准确正则性. 它的证明远远超出本书范围, 例如可见 Grisvard [42].

\setcounter{statement}{7}
\renewcommand{\theHstatement}{theorem.\thesection.\arabic{statement}}
\begin{Theorem}\label{thm:laplace-dirichlet-regularity}
$1^\circ)$ 设 $\Omega$ 是 $\mathbb R^N$ 的有界开子集, 其边界 $\Gamma$ 对某个整数 $k\geq0$ 属于 $\mathcal C^{k+1,1}$ 类. 假设 Problem~\eqref{eq:laplace-dirichlet-pde}--\eqref{eq:laplace-dirichlet-boundary} 的数据 $f$ 和 $g$ 满足

\[
f\in W^{k,p}(\Omega),\qquad g\in W^{k+2-1/p,p}(\Gamma),
\]

其中实数 $p$ 满足 $1<p<\infty$. 则 $u\in W^{k+2,p}(\Omega)$, 并且存在常数 $C=C(k,p,\Omega)$, 使得

\begin{equation}\label{eq:dirichlet-regularity-wkp}
\|u\|_{k+2,p,\Omega}\leq C\{\|f\|_{k,p,\Omega}+\|g\|_{k+2-1/p,p,\Gamma}\}.
\end{equation}

$2^\circ)$ 当 $\Omega$ 是没有凹角的有界二维多边形时, 存在依赖于 $\Gamma$ 最大内角的实数 $p_\Omega>2$, 使得只要 $f\in L^p(\Omega)$, 并且 $(g|_{\Gamma_j};\ 1\leq j\leq J)\in\prod_{j=1}^J W^{2-1/p,p}(\Gamma_j)$ 满足 Remark~\ref{rem:polygon-boundary-traces} 的匹配条件, 就有

\[
u\in W^{2,p}(\Omega),\qquad1<p<p_\Omega.
\]

$3^\circ)$ 如果 $\Omega$ 是三维有界凸多面体, 则 $2^\circ)$ 的结论对于齐次 Dirichlet 问题 $(g=0)$ 仍然成立.
\end{Theorem}

\MathBookSourcePage{27}

\setcounter{statement}{1}
\renewcommand{\theHstatement}{remark.\thesection.\arabic{statement}}
\begin{Remark}\label{rem:laplace-isomorphism-convex-domain}
令 $g=0$ 应用 Theorem~\ref{thm:laplace-dirichlet-regularity}, 立即可知: 当 $\Omega$ 是 $\mathbb R^2$ 中的有界凸多边形(或 $\mathbb R^3$ 中的多面体)时, 映射 $u\mapsto\Delta u$ 是从 $W^{2,p}(\Omega)\cap W_0^{1,p}(\Omega)$ 到 $L^p(\Omega)$ 的同构, 其中对某个 $\varepsilon>0$, $p\in(1,2+\varepsilon]$. 当 $\Omega$ 的边界属于 $\mathcal C^{1,1}$ 类时, 这个同构对所有 $p\in(1,\infty)$ 都成立.
\end{Remark}

\subsection{Example 2: Laplace 算子的 Neumann 问题}\label{subsec:laplace-neumann-example}

这里进一步假设 $\Omega$ 连通, 并处理如下非齐次 Neumann 问题:

求 $u$, 使得

\noindent Problem (N):
\begin{align}
-\Delta u&=f&&\text{in }\Omega,\label{eq:laplace-neumann-pde}\\
\partial u/\partial n&=g&&\text{on }\Gamma,\label{eq:laplace-neumann-boundary}\\
\int_\Omega f\,dx+\langle g,1\rangle_\Gamma&=0,
&& f\in L^2(\Omega),\ g\in H^{-1/2}(\Gamma).\label{eq:laplace-neumann-compatibility}
\end{align}

由于 Problem (N) 只涉及 $u$ 的导数, 它的解显然不可能唯一. 为绕开这一困难, 在商空间 $H^1(\Omega)/\mathbb R$ 中求解 $u$, 并为该空间赋予商范数

\begin{equation}\label{eq:h1-quotient-norm}
\|\dot v\|_{H^1(\Omega)/\mathbb R}=\inf_{v\in\dot v}\|v\|_{1,\Omega}.
\end{equation}

下面的 Theorem 给出该空间的一个重要性质; 其证明可见 Nečas [58].

\setcounter{statement}{8}
\renewcommand{\theHstatement}{theorem.\thesection.\arabic{statement}}
\begin{Theorem}\label{thm:h1-quotient-hilbert}
设 $\Omega$ 是 $\mathbb R^N$ 的有界, 连通且 Lipschitz 连续的开子集. 对于商范数
\eqref{eq:h1-quotient-norm}, 空间 $H^1(\Omega)/\mathbb R$ 是 Hilbert 空间. 此外, 在这个空间上,
泛函 $\dot v\mapsto|v|_{1,\Omega}$ 是与 \eqref{eq:h1-quotient-norm} 等价的范数.
\end{Theorem}

利用这个空间, 可以把 Problem (N) 化为 Problem~\eqref{eq:abstract-variational} 的抽象形式.
令 $V=H^1(\Omega)/\mathbb R$,

\[
a(\dot u,\dot v)=(\mathbf{grad}\,u,\mathbf{grad}\,v),
\]

并令

\begin{equation}\label{eq:neumann-quotient-functional}
l:\dot v\mapsto(f,v)+\langle g,v\rangle_\Gamma\qquad\forall v\in\dot v.
\end{equation}

注意, 由于相容条件 \eqref{eq:laplace-neumann-compatibility}, 式
\eqref{eq:neumann-quotient-functional} 的右端与所取的具体代表元 $v\in\dot v$ 无关.
此外, 由 \eqref{eq:trace-space-norm} 有

\[
|(f,v)+\langle g,v\rangle_\Gamma|\leq(\|f\|_{0,\Omega}+\|g\|_{-1/2,\Gamma})\inf_{v\in\dot v}\|v\|_{1,\Omega},
\]

因此 $l\in V'$.

\MathBookSourcePage{28}

从而

\begin{equation}\label{eq:neumann-functional-bound}
\|l\|_{V'}\leq\|f\|_{0,\Omega}+\|g\|_{-1/2,\Gamma}.
\end{equation}

显然, $a(\dot u,\dot v)$ 在 $V\times V$ 上连续, 并且由
Theorem~\ref{thm:h1-quotient-hilbert},

\[
a(\dot v,\dot v)=|v|_{1,\Omega}^2\geq C_1\|\dot v\|_{H^1(\Omega)/\mathbb R}^2.
\]

因此, 由 Lax \& Milgram Theorem, 如下问题

\begin{equation}\label{eq:neumann-variational-problem}
(N')\quad\text{Find }\dot u\in H^1(\Omega)/\mathbb R\text{ satisfying}\qquad a(\dot u,\dot v)=\langle l,\dot v\rangle\quad\forall\dot v\in H^1(\Omega)/\mathbb R
\end{equation}

在 $H^1(\Omega)/\mathbb R$ 中存在唯一解 $\dot u$.

下面解释 Problem $(N')$. 把 $v$ 限制在 $\mathcal D(\Omega)$ 中时,
式 \eqref{eq:neumann-variational-problem} 给出 \eqref{eq:laplace-neumann-pde}.
接着, 将式 \eqref{eq:laplace-neumann-pde} 与 $v$ 作标量积, 并与
\eqref{eq:neumann-variational-problem} 比较, 得到

\begin{equation}\label{eq:neumann-green-weak-boundary}
(\mathbf{grad}\,u,\mathbf{grad}\,v)=-(\Delta u,v)+\langle g,v\rangle_\Gamma\qquad\forall v\in H^1(\Omega).
\end{equation}

因此, Problem $(N')$ 等价于在 $H^1(\Omega)$ 中求满足
\eqref{eq:laplace-neumann-pde} 和 \eqref{eq:neumann-green-weak-boundary} 的 $u$.

还需把 \eqref{eq:neumann-green-weak-boundary} 解释为边界条件. 在这一阶段, 如果不假设
$u\in H^2(\Omega)$, 就无法做到这一点. 此时 Green 公式
\eqref{eq:green-first-identity} 给出

\[
\int_\Gamma(\partial u/\partial n)v\,ds=\langle g,v\rangle_\Gamma\qquad\forall v\in H^1(\Omega),
\]

即 $\partial u/\partial n=g$ on $\Gamma$. 因此, Problems (N) 和 $(N')$ 等价. 当然, 这并不完全令人满意, 因为 Problem (N) 的解是否存在取决于 $(N')$ 的解的正则性. 虽然这种正则性通常确实成立, 但下一段中更强有力的工具将消除这个额外的光滑性假设.

现在考察 Problem $(N')$ 的解 $\dot u$ 对数据的依赖性. 根据 Lax--Milgram
Theorem~\ref{thm:lax-milgram}, 式 \eqref{eq:neumann-functional-bound} 以及
Theorem~\ref{thm:h1-quotient-hilbert} 的范数等价性, 得到

\[
|u|_{1,\Omega}\leq C_2(\|f\|_{0,\Omega}+\|g\|_{-1/2,\Gamma}).
\]

于是证明了如下结果.

\setcounter{statement}{1}
\renewcommand{\theHstatement}{proposition.\thesection.\arabic{statement}}
\begin{Proposition}\label{prop:laplace-neumann-wellposedness}
Problem $(N')$ 在 $H^1(\Omega)/\mathbb R$ 中存在唯一解 $\dot u$, 并且这个解连续依赖于数据:

\MathBookSourcePage{29}

\begin{equation}\label{eq:neumann-stability-h1}
|u|_{1,\Omega}\leq C(\|f\|_{0,\Omega}+\|g\|_{-1/2,\Gamma})\qquad\forall u\in\dot u.
\end{equation}

此外, 当 $\dot u\in H^2(\Omega)/\mathbb R$ 时, 它也是 Problem (N) 的唯一解.
\end{Proposition}

与 Dirichlet 问题类似, 当 Problem $(N')$ 的数据具有额外光滑性时, 它的解具有更高正则性.
由 Grisvard [42] 给出的准确结果与 Theorem~\ref{thm:laplace-dirichlet-regularity} 十分相似.

\setcounter{statement}{9}
\renewcommand{\theHstatement}{theorem.\thesection.\arabic{statement}}
\begin{Theorem}\label{thm:laplace-neumann-regularity}
$1^\circ)$ 设 $\Omega$ 如 Theorem~\ref{thm:laplace-dirichlet-regularity} 中所述, 并假设
Problem~\eqref{eq:neumann-variational-problem} 的数据 $f$ 和 $g$ 满足

\[
f\in W^{k,p}(\Omega),\qquad g\in W^{k+1-1/p,p}(\Gamma),\qquad1<p<\infty.
\]

则 $\dot u\in W^{k+2,p}(\Omega)/\mathbb R$, 并且存在常数 $C=C(k,p,\Omega)$, 使得

\begin{equation}\label{eq:neumann-regularity-wkp}
\|\dot u\|_{W^{k+2,p}(\Omega)/\mathbb R}\leq C\{\|f\|_{k,p,\Omega}+\|g\|_{k+1-1/p,p,\Gamma}\}.
\end{equation}

$2^\circ)$ 当 $\Omega$ 是没有凹角的有界二维多边形时, 存在依赖于 $\Gamma$ 最大内角的实数 $p_\Omega>2$, 使得只要 $f\in L^p(\Omega)$, 并且 $(g|_{\Gamma_j};\ 1\leq j\leq J)\in\prod_{j=1}^J W^{1-1/p,p}(\Gamma_j)$, 就有

\[
\dot u\in W^{2,p}(\Omega)/\mathbb R,\qquad1<p<p_\Omega.
\]

$3^\circ)$ 如果 $\Omega$ 是 $\mathbb R^3$ 中的有界凸多面体, 则 $2^\circ)$ 的结论对齐次 Neumann 问题 $(g=0)$ 成立.
\end{Theorem}

\subsection{Example 3: 双调和算子的 Dirichlet 问题}\label{subsec:biharmonic-dirichlet-example}

考虑如下非齐次问题:

给定 $H^{-2}(\Omega)$ 中的 $f$, $H^{3/2}(\Gamma)$ 中的 $g_1$ 以及 $H^{1/2}(\Gamma)$ 中的 $g_2$, 求 $u$, 使得

\noindent Problem (B):
\begin{align}
\Delta^2u&=f&&\text{in }\Omega,\label{eq:biharmonic-equation}\\
u&=g_1&&\text{on }\Gamma,\label{eq:biharmonic-dirichlet-boundary}\\
\partial u/\partial n&=g_2&&\text{on }\Gamma.\label{eq:biharmonic-normal-boundary}
\end{align}

与这个问题自然对应的函数空间是 $H_0^2(\Omega)$, 双线性形式为

\[
a(u,v)=(\Delta u,\Delta v).
\]

该形式在 $H_0^2(\Omega)$ 上是椭圆的, 因为映射 $v\mapsto\|\Delta v\|_{0,\Omega}$ 是 $H_0^2(\Omega)$ 上与范数 $\|\cdot\|_{2,\Omega}$ 等价的范数. 事实上, 对 $\mathcal D(\Omega)$ 中的 $v$, 通过分部积分并交换导数, 很容易证明

\begin{equation}\label{eq:biharmonic-laplacian-norm}
\|\Delta v\|_{0,\Omega}^2=|v|_{2,\Omega}^2.
\end{equation}

由稠密性, 同样的结果对 $H_0^2(\Omega)$ 中的函数成立. 范数等价性来自
Poincaré Theorem~\ref{thm:poincare-friedrichs}.

\MathBookSourcePage{30}

根据 Theorem~\ref{thm:normal-derivative-trace}, 如果 $\Gamma$ 属于 $\mathcal C^{1,1}$ 类,
则存在函数 $u_0\in H^2(\Omega)$, 使得

\begin{equation}\label{eq:biharmonic-lifting-data}
u_0=g_1\quad\text{on }\Gamma,\qquad\partial u_0/\partial n=g_2\quad\text{on }\Gamma.
\end{equation}

于是转而考虑如下问题:

\noindent Problem $(B')$: 求 $u\in H^2(\Omega)$, 使得
\begin{align}
u-u_0&\in H_0^2(\Omega),\label{eq:biharmonic-lifted-space}\\
a(u-u_0,v)&=\langle f,v\rangle-a(u_0,v)
&&\forall v\in H_0^2(\Omega).\label{eq:biharmonic-variational}
\end{align}

由 Lax--Milgram Theorem~\ref{thm:lax-milgram}, Problem $(B')$ 在 $H^2(\Omega)$ 中恰有一个解 $u$.
由 \eqref{eq:biharmonic-lifting-data} 和 \eqref{eq:biharmonic-lifted-space}, $u$ 满足边界条件
$u=g_1$ 和 $\partial u/\partial n=g_2$ on $\Gamma$. 此外, 把
\eqref{eq:biharmonic-variational} 的测试函数限制到 $\mathcal D(\Omega)$, 得到
$\Delta^2u=f$ in $H^{-2}(\Omega)$. 因此 $u$ 是 (B) 的解.

反过来, 与 Laplace 算子的情形一样, 可以证明 Problem (B) 在 $H^2(\Omega)$ 中至多有一个解.
由 \eqref{eq:biharmonic-variational} 和范数等价性得到估计

\[
\|u\|_{2,\Omega}\leq C_1(\|f\|_{-2,\Omega}+\|u_0\|_{2,\Omega})
\qquad\forall u_0\text{ satisfying }\eqref{eq:biharmonic-lifting-data},
\]

即

\[
\|u\|_{2,\Omega}\leq C_2(\|f\|_{-2,\Omega}+\|g_1\|_{3/2,\Gamma}+\|g_2\|_{1/2,\Gamma}).
\]

这些结果概括为下面的 Proposition.

\setcounter{statement}{2}
\renewcommand{\theHstatement}{proposition.\thesection.\arabic{statement}}
\begin{Proposition}\label{prop:biharmonic-dirichlet-wellposedness}
如果 $\Gamma$ 属于 $\mathcal C^{1,1}$ 类, 则 Problem (B) 在 $H^2(\Omega)$ 中恰有一个解 $u$, 并且满足如下估计:

\begin{equation}\label{eq:biharmonic-stability-h2}
\|u\|_{2,\Omega}\leq C(\|f\|_{-2,\Omega}+\|g_1\|_{3/2,\Gamma}+\|g_2\|_{1/2,\Gamma}).
\end{equation}
\end{Proposition}

上述分析假设 $\Gamma$ 属于 $\mathcal C^{1,1}$ 类, 因而不允许有角点. 这个假设在提升算子 $(g_1,g_2)\mapsto u_0$ 中起关键作用. 当然, 如果这些非齐次数据直接以函数 $u_0\in H^2(\Omega)$ 的形式给出, 且

\[
u=u_0\quad\text{on }\Gamma\quad\text{and}\quad\partial u/\partial n=\partial u_0/\partial n\quad\text{on }\Gamma,
\]

那么 Proposition~\ref{prop:biharmonic-dirichlet-wellposedness} 也适用于 Lipschitz 连续定义域,
其中用 $u_0$ 代替 $g_1$ 和 $g_2$. 否则, 必须略微修改 Problem (B) 的陈述.
假设 $\Omega$ 是有界二维多边形. 按照 Remark~\ref{rem:polygon-boundary-traces},
以 $\Gamma_j$ 和 $S_j$, $1\leq j\leq J$, 分别表示 $\Gamma$ 的边和顶点. 取 $J$ 个函数:

\MathBookSourcePage{31}

\[
h_j\in H^{3/2}(\Gamma_j),\qquad1\leq j\leq J,
\]

它们满足匹配条件 $h_j(S_{j+1})=h_{j+1}(S_{j+1})$, 再取 $J$ 个函数

\[
g_j\in H^{1/2}(\Gamma_j),\qquad1\leq j\leq J,
\]

并考虑如下问题:

\noindent Problem $(B'')$ 使用 \eqref{eq:biharmonic-equation}, 并满足
\begin{align}
u&=h_j&&\text{on }\Gamma_j,\quad1\leq j\leq J,
\label{eq:biharmonic-polygon-dirichlet}\\
\partial u/\partial n&=g_j&&\text{on }\Gamma_j,\quad1\leq j\leq J.
\label{eq:biharmonic-polygon-normal}
\end{align}

由 Remark~\ref{rem:polygon-boundary-traces}, 存在函数 $u_0\in H^2(\Omega)$, 使得

\[
u_0=h_j,\qquad\partial u_0/\partial n=g_j\quad\text{on }\Gamma_j,\qquad1\leq j\leq J,
\]

并且

\[
\|u_0\|_{2,\Omega}\leq C_3\left\{\sum_{j=1}^J(\|h_j\|_{3/2,\Gamma_j}^2+\|g_j\|_{1/2,\Gamma_j}^2)\right\}^{1/2}.
\]

因此, Proposition~\ref{prop:biharmonic-dirichlet-wellposedness} 的结论(用函数 $h_j$ 和 $g_j$
代替 $g_1$ 和 $g_2$)在这种情形下也适用于 Problem $(B'')$.

当 $\Omega$ 的边界充分光滑时, 可以得到关于 $u$ 的正则性的更多信息.

\setcounter{statement}{10}
\renewcommand{\theHstatement}{theorem.\thesection.\arabic{statement}}
\begin{Theorem}\label{thm:biharmonic-smooth-boundary-regularity}
设 $\Omega$ 是 $\mathbb R^N$ 的有界开子集, 其边界 $\Gamma$ 对某个整数 $k\geq-2$ 属于 $\mathcal C^{k+3,1}$ 类, 并假设双调和 Problem (B) 的数据 $f,g_1,g_2$ 满足

\[
f\in H^k(\Omega),\qquad g_1\in H^{k+7/2}(\Gamma),\qquad g_2\in H^{k+5/2}(\Gamma).
\]

则 $u\in H^{k+4}(\Omega)$, 并且存在常数 $C=C(k,\Omega)$, 使得

\begin{equation}\label{eq:biharmonic-regularity-hk}
\|u\|_{k+4,\Omega}\leq C\{\|f\|_{k,\Omega}+\|g_1\|_{k+7/2,\Gamma}+\|g_2\|_{k+5/2,\Gamma}\}.
\end{equation}
\end{Theorem}

即使 $\Gamma$ 有角点, 对于具有齐次边界条件的双调和问题,
Proposition~\ref{prop:biharmonic-dirichlet-wellposedness} 的结论仍可加强.

\setcounter{statement}{11}
\renewcommand{\theHstatement}{theorem.\thesection.\arabic{statement}}
\begin{Theorem}\label{thm:biharmonic-convex-polygon-regularity}
假设 $\Omega$ 是没有凹角的有界二维多边形. 则映射 $u\mapsto\Delta^2u$ 是从
$H^3(\Omega)\cap H_0^2(\Omega)$ 到 $H^{-1}(\Omega)$ 的同构.
\end{Theorem}

最后这个结果对于建立平面凸多边形上 Stokes 问题解的正则性是基本的(参见 Grisvard [43]).
